\section{Critical radius content and an elliptic residue curve}\label{sec:critical-content}

The common content of the four integral radii produces a finite-field
curve independent of the original denominator. Its elliptic structure
gives a classification at every content prime. A second elliptic curve
records the square class of the opposite-product parameter.

\subsection{Critical reduction}

Write the primitive integral equations as
\begin{equation}\label{crit:integral}
\begin{aligned}
 A_0^2&=X^2+Y^2,& A_1^2&=(X-L)^2+Y^2,\\
 A_2^2&=X^2+(Y-L)^2,& A_3^2&=(X-L)^2+(Y-L)^2,
\end{aligned}
\end{equation}
with corner order \(00,10,01,11\), positive radii, \(0<X,Y<L\),
and \(\gcd(X,Y,L)=1\). Thus \(r_j=A_j/L\). Put
\[
 G=\gcd_j A_j,\qquad H=L/G^2,\qquad a_j=A_j/G.
\]
Recall from Proposition~\ref{int:content} that \(G^2\mid L\) and
\(p\mid G\Rightarrow p\equiv1\pmod4\). The local argument also
fixes the valuations needed here. If an odd \(p\mid G\) did not
divide \(L\), subtraction would give \(X=Y=L/2\pmod p\), contrary
to \(X^2+Y^2=0\). Hence \(p\mid L\), and primitivity makes
\(X,Y\) units. Consequently \((Y/X)^2=-1\pmod p\) and
\[
 v_p(A_1^2-A_0^2)=v_p\bigl(L(L-2X)\bigr)=v_p(L)\ge2v_p(G).
\]
The prime 2 cannot divide \(G\), since the equations modulo four
would then force \(X,Y,L\) all even.

\begin{definition}
A prime \(p\mid G\) is \emph{critical} if \(v_p(L)=2v_p(G)\).
\end{definition}

At a critical prime set \(s=v_p(G)>0\) and, in \(\F_p\),
\[
 \iota=Y/X,\qquad\iota^2=-1,\qquad b=-2HX\ne0,\qquad c=a_0^2.
\]
The four normalized radii satisfy
\begin{equation}\label{crit:pattern}
 (a_0^2,a_1^2,a_2^2,a_3^2)
 =\bigl(c,c+b,c+\iota b,c+(1+\iota)b\bigr).
\end{equation}
The shifts are distinct, so at most one normalized radius vanishes.
At a noncritical content prime, \(p\mid H\); all \(a_j^2\) are then
equal modulo \(p\), and \(\gcd_j a_j=1\) makes every \(a_j\) a unit.

\subsection{The fixed elliptic curve}

Over a field \(k\) of characteristic different from two containing
\(\iota^2=-1\), define
\begin{equation}\label{crit:curve}
 \mathcal C_\iota:\quad
 \begin{cases}
 V_2^2-V_0^2=\iota(V_1^2-V_0^2),\\
 V_3^2-V_0^2=(1+\iota)(V_1^2-V_0^2)
 \end{cases}
 \quad\subset\PP^3_k.
\end{equation}
Write \(b=V_1^2-V_0^2\), and let \(U_\iota\) be the open subset
\(bV_0V_1V_2V_3\ne0\).

\begin{proposition}\label{prop:crit-elliptic}
The curve \(\mathcal C_\iota\) is smooth, geometrically connected,
and \(k\)-isomorphic to the smooth projective curve
\(E:y^2=x^3-x\). Its eight points with \(b=0\) are
\([1:\epsilon_1:\epsilon_2:\epsilon_3]\), with \(\epsilon_j=\pm1\).
\end{proposition}

\begin{proof}
The parameter \(T=V_0^2/b\) expresses its function field as
\[
 k(T)\left(\sqrt{\frac{T+1}{T}},\sqrt{\frac{T+\iota}{T}},
                  \sqrt{\frac{T+1+\iota}{T}}\right).
\]
The three square classes are geometrically independent, by their
respective odd valuations at \(-1,-\iota,-1-\iota\).
The degree-eight cover has four branch values
\(0,-1,-\iota,-1-\iota\), each with four points of ramification
index two, and is unramified at infinity. Riemann--Hurwitz gives
\(2g-2=-16+16=0\).
The diagonal pencil of the two quadrics has coefficients
\[
 (\iota-1)\lambda+\iota\mu,\quad
 -\iota\lambda-(1+\iota)\mu,\quad\lambda,\quad\mu.
\]
Their zero positions are distinct. A singular point would therefore
have at most one nonzero coordinate, and no such point satisfies the
two equations. Thus the displayed intersection is smooth.

Even coordinate-sign changes, modulo simultaneous sign change, give
a degree-four quotient
\begin{equation}\label{crit:quartic}
 W^2=T(T+1)(T+\iota)(T+1+\iota),\qquad
 W=V_0V_1V_2V_3/b^2.
\end{equation}
A nonidentity element has a projective fixed point only if two
coordinates vanish, which is impossible. After choosing a rational
origin, each such involution is translation by nonzero two-torsion.
Indeed an automorphism \(P\mapsto\alpha(P)+Q\) with
\(\alpha\ne1\) has a fixed point because \(1-\alpha\) is surjective.
The quotient group is therefore the entire two-torsion translation
group. Factoring multiplication by two shows that the quotient is
\(k\)-isomorphic to \(\mathcal C_\iota\).

In \eqref{crit:quartic}, set \(R=(1+\iota)T+\iota\), \(S=2W\).
Then \(S^2=1-R^4\), and
\[
 X_E=\frac{2(1+S)}{R^2},\quad Y_E=\frac{4(1+S)}{R^3},\qquad
 R=\frac{2X_E}{Y_E},\quad S=\frac{X_E^2-4}{X_E^2+4}
\]
identify its function field with \(Y_E^2=X_E^3+4X_E\).
Scaling by \(x=X_E/(1+\iota)^2\), \(y=Y_E/(1+\iota)^3\), and
using \((1+\iota)^4=-4\), gives \(y^2=x^3-x\).
The identifications extend to the smooth projective curves.
\end{proof}

For \(p\equiv1\pmod4\), write \(\chi_p\) for the quadratic
character, with \(\chi_p(0)=0\), and put
\[
 \delta_p=\begin{cases}
 1,&\chi_p(\iota)=\chi_p(1+\iota)=1,\\0,&\text{otherwise}.
 \end{cases}
\]
The zero-coordinate fibres are disjoint; each has four points when
\(\delta_p=1\), and none otherwise. For example \(V_0=0,V_1=1\)
gives \(V_2^2=\iota,V_3^2=1+\iota\). Dividing by a nonzero
coordinate in each other fibre gives the same criterion. Therefore
\begin{equation}\label{crit:count}
 \#U_\iota(\F_p)=\#E(\F_p)-8-16\delta_p.
\end{equation}
Each admissible \(T=c/b\) gives eight projective points. Equivalently,
\((\#E(\F_p)-8-16\delta_p)/8\) counts the \(T\) for which
\(T,T+1,T+\iota,T+1+\iota\) are nonzero and have a common quadratic
character. The square class of \(b\) supplies either common character.

\subsection{The complete prime classification}

\begin{theorem}\label{thm:crit-basic}
For a prime \(p\equiv1\pmod4\), the following are equivalent:
\begin{enumerate}
\item \(p\notin\{5,13\}\).
\item A pattern \eqref{crit:pattern} exists with \(b\ne0\) and all
      four entries nonzero squares in \(\F_p\).
\item For every \(s\ge1\), there are infinitely many primitive
      triples \((X,Y,L)\), with \(0<X,Y<L\) and \(v_p(L)=2s\),
      such that all four radicands in \eqref{crit:integral} have
      roots in \(\Q_p\) of valuation exactly \(s\).
\end{enumerate}
At 5 and 13 no critical pattern exists even with zero entries.
In particular, for a rational four-distance point,
\[
 p\in\{5,13\},\quad p\mid G\quad\Longrightarrow\quad p\mid H.
\]
\end{theorem}

\begin{proof}
Hasse's bound \(\#E(\F_p)\ge p+1-2\sqrt p\)
\cite{silverman2009} exceeds 24 for \(p\ge37\).
Equation~\eqref{crit:count} then gives a point in \(U_\iota\).
The remaining existence cases are
\[
\begin{array}{c|c|c|c|c}
p&\iota&(V_0,V_1,V_2,V_3)&b&c\\\hline
17&4&(1,2,8,4)&3&1\\
29&12&(2,8,12,1)&2&4.
\end{array}
\]
Interchanging \(V_1,V_2\) replaces \(\iota\) by
\(\iota^{-1}=-\iota\), covering the other root of minus one.

For the exclusions, division by a nonzero square reduces \(b\) to
one of its two square classes. The square sets including zero are
\[
 S_5=\{0,1,4\},\qquad S_{13}=\{0,1,3,4,9,10,12\}.
\]
All possibilities after the first two square tests are listed below:
\[
\begin{array}{c|c|c|l}
p&\iota&b&c\ ;\ (c+\iota b,c+(1+\iota)b)\\\hline
5&2&1&0;(2,3)\quad4;(1,2)\\
5&2&2&4;(3,0)\\
13&5&1&0;(5,6)\quad3;(8,9)\quad9;(1,2)\quad12;(4,5)\\
13&5&2&1;(11,0)\quad10;(7,9)\quad12;(9,11).
\end{array}
\]
Every pair contains a nonsquare.

To realize an allowed signed unit pattern \((V_0,V_1,V_2,V_3)\),
put \(c=V_0^2,b=V_1^2-V_0^2\), and fix \(s\ge1\).
Choose positive integers \(X,Y\) with
\begin{equation}\label{crit:lift-coordinates}
 X\equiv-b/2\pmod p,\qquad
 Y\equiv\widetilde\iota X+p^{2s}z\pmod{p^{2s+1}},\qquad
 z\equiv\frac{c}{2\iota X}\pmod p,
\end{equation}
where Hensel's lemma supplies
\(\widetilde\iota^2\equiv-1\pmod{p^{2s+1}}\).
For sufficiently large positive integers \(k\), put
\begin{equation}\label{crit:lift-denominator}
 M=1+pXYk,\qquad L=p^{2s}M.
\end{equation}
Then \(L>X,Y\), \(v_p(L)=2s\), and \(\gcd(X,Y,L)=1\), since
\(p\nmid XY\) and \(M\equiv1\pmod{XY}\).
Expanding \eqref{crit:lift-coordinates} gives
\((X^2+Y^2)/p^{2s}\equiv c\pmod p\).
Subtracting \(2LX\) or \(2LY\), and adding \(L^2\), gives the
other three reductions in \eqref{crit:pattern}.
Hensel's lemma now supplies roots \(A_j=p^sv_j\), with
\(v_j\in\Z_p^\times\) and \(v_j\equiv V_j\pmod p\), of the four
original radicands. Varying \(k\) gives infinitely many coordinate
pairs. Conversely, roots with these valuations reduce to a critical
pattern, to which the exclusions at 5 and 13 apply.
\end{proof}

\begin{corollary}\label{cor:crit-crt}
Let \(\mathcal P\) be a finite set of primes congruent to one modulo
four, excluding 5 and 13, and prescribe \(s_p\ge1\).
There are infinitely many primitive triples \((X,Y,L)\), with
\(0<X,Y<L\), for which \(v_p(L)=2s_p\) and all four original
radicands have roots of valuation \(s_p\) in \(\Q_p\), simultaneously
for every \(p\in\mathcal P\).
\end{corollary}

\begin{proof}
Choose \(X,Y\) by the Chinese remainder theorem using
\eqref{crit:lift-coordinates} at the selected primes. If
\(D=\gcd(X,Y)\), no selected prime divides \(D\). The congruences
\[
 L\equiv p^{2s_p}\pmod{p^{2s_p+1}}\quad(p\in\mathcal P),
 \qquad L\equiv1\pmod D
\]
are compatible. Arbitrarily large solutions give the strict
inequalities and primitivity. The prescribed unit-square reductions
persist, so Hensel's lemma supplies the roots.
\end{proof}

\subsection{The opposite product and the prime 17}

Define
\begin{equation}\label{crit:product}
 t=r_0r_3+r_1r_2=\frac{Q}{G^2H^2},\qquad
 Q=a_0a_3+a_1a_2,\quad Q'=a_0a_3-a_1a_2.
\end{equation}
At a critical prime, multiplication in \(\F_p\) gives
\begin{equation}\label{crit:product-unit}
 QQ'=c\bigl(c+(1+\iota)b\bigr)-(c+b)(c+\iota b)
     =-\iota b^2\ne0.
\end{equation}
Thus both products are units, even with a zero radius reduction.
For \(s=v_p(G)\), the denominator unit in \eqref{crit:product} is a
square, and hence
\begin{equation}\label{crit:product-class}
 v_p(t)=-2s,\qquad [p^{2s}t]=[Q]
 \text{ in }\Q_p^\times/\Q_p^{\times2},\qquad
 \frac{t+1}{t}\in\Q_p^{\times2}.
\end{equation}
The last assertion follows from
\((t+1)/t\in1+p^{2s}\Z_p\).

\begin{theorem}\label{thm:crit-17}
If 17 is a critical content prime of a rational four-distance point,
then none of \(t,2t,t+1,2(t+1)\) is a square in
\(\Q_{17}\). Therefore, if any one is a rational square,
\(17\mid G\) implies \(17\mid H\).
At a critical 17, with \(s=v_{17}(G)\), one also has
\[
 a_0^2+a_3^2\equiv a_1^2+a_2^2\equiv0\pmod{17},\qquad
 17^{2s+1}\mid(2X-L)^2+(2Y-L)^2.
\]
\end{theorem}

\begin{proof}
Take \(\iota=4\). The square set is
\(S_{17}=\{0,1,2,4,8,9,13,15,16\}\), and the scaling classes of
\(b\) are represented by 1 and 3. Successive square tests give
\[
\begin{array}{c|c|c|c}
b&c,c+b\in S_{17}&\text{also }c+4b\in S_{17}
 &\text{also }c+5b\in S_{17}\\\hline
1&\{0,1,8,15,16\}&\{0,15\}&\varnothing\\
3&\{1,13,15,16\}&\{1,13\}&\{1\}.
\end{array}
\]
Thus all coordinates are nonzero, and division by \(a_0\) gives
\([1:\pm2:\pm8:\pm4]\). Hence
\(Q/a_0^2\in\{3,5,12,14\}\), all nonsquares modulo 17.
Interchanging the middle coordinates handles \(\iota=-4\).
Since 2 and \((t+1)/t\) are local squares,
\eqref{crit:product-class} gives the four exclusions.

The same table yields \(c/b=-(1+\iota)/2\), so
\(c=H(X+Y)\) and \(2c+(1+\iota)b=0\).
The identity
\begin{equation}\label{crit:center-identity}
 \frac{(2X-L)^2+(2Y-L)^2}{G^2}
 =4a_0^2-4H(X+Y)+2G^2H^2
\end{equation}
proves the divisibility. At a noncritical 17, its right side is a unit,
since \(17\mid H\) and \(a_0\) is a unit, so the corresponding
valuation is exactly \(2s\).
\end{proof}

\subsection{The normalized product-square cover}

On \(\mathcal C_\iota\) put
\(f=V_0V_3+V_1V_2\), \(f'=V_0V_3-V_1V_2\), so that
\(ff'=-\iota b^2\). Let \(\widetilde{\mathcal C}_\iota\) be the
smooth projective curve obtained by adjoining
\begin{equation}\label{crit:double}
 S^2=2f/V_0^2.
\end{equation}
On \(V_j\ne0\), the same extension is \(S_j^2=2f/V_j^2\).

\begin{proposition}\label{prop:crit-double}
The map \(\widetilde{\mathcal C}_\iota\to\mathcal C_\iota\) is a
geometrically connected unramified double cover of genus-one curves
with rational points. For \(p\equiv1\pmod4\), its number of
rational points above \(b=0\) is
\[
 B_p=\begin{cases}16,&p\equiv1\pmod8,\\8,&p\equiv5\pmod8.
 \end{cases}
\]
Moreover,
\begin{equation}\label{crit:double-count}
 \#\{V\in\mathcal C_\iota(\F_p):b(V)\ne0,\ \chi_p(2f(V))=1\}
 =\frac{\#E(\F_p)-B_p}{2}.
\end{equation}
The character is unchanged by rescaling projective coordinates.
\end{proposition}

\begin{proof}
At \(b=0\), the four zeros of \(f\) are
\([1:\epsilon_1:\epsilon_2:-\epsilon_1\epsilon_2]\).
There \(f'\ne0\), and \(b/V_0^2\) is a uniformizer, because the
map \(T=V_0^2/b\) is unramified at infinity. Thus \(ff'=-\iota b^2\)
makes every zero of \(f/V_0^2\) have order two. On \(V_0=0\), one has \(f=V_1V_2\ne0\), so its poles
also have even order, proving that the cover is
geometrically unramified.

If the extension were geometrically trivial, its square root times
\(V_0\) would be a regular global section of \(\OO_{\mathcal C_\iota}(1)\)
vanishing precisely at the four displayed points. This degree-four
line bundle has four sections by Riemann--Roch; the four independent
coordinate sections span them. But the coordinate matrix of these
points has determinant \(\pm16\), so no linear form vanishes at all
four. The cover is connected. Riemann--Hurwitz gives genus one, and
\(S=\pm2\) above \([1:1:1:1]\) gives rational points.
With compatible origins the map is a degree-two isogeny.

At the four points with \(b=0,f\ne0\), the values
\(2f/V_0^2=\pm4\) give eight rational preimages. At each remaining
point, the normalized local equation has residue
\[
 \left(\frac{S}{b/V_0^2}\right)^2
 =\frac{-2\iota V_0^2}{f'}=\pm\iota.
\]
It gives two rational points precisely when \(\iota\) is square,
equivalently \(p\equiv1\pmod8\). This proves the boundary count.
Isogenous elliptic curves over a finite field have equal point
counts. Apply this to the double cover and
Proposition~\ref{prop:crit-elliptic}. For \(b\ne0\), one has
\(f\ne0\), and each square value has exactly two preimages, proving
\eqref{crit:double-count}.
\end{proof}

\begin{theorem}\label{thm:crit-refined}
For a prime \(p\equiv1\pmod4\), the following are equivalent:
\begin{enumerate}
\item \(p\notin\{5,13,17\}\).
\item There is a signed critical pattern with every \(V_j\ne0\)
      and \(2(V_0V_3+V_1V_2)\in\F_p^{\times2}\).
\item For every \(s\ge1\), there are infinitely many primitive
      strictly interior rational coordinate pairs \((X/L,Y/L)\),
      with \(v_p(L)=2s\), and roots \(A_j\in\Q_p\) of
      \eqref{crit:integral}, each of valuation \(s\), such that
      \(t=(A_0A_3+A_1A_2)/L^2\) satisfies
      \(2t,2(t+1)\in\Q_p^{\times2}\).
\end{enumerate}
\end{theorem}

\begin{proof}
If \(p\equiv5\pmod8\), zero-coordinate base points are impossible
because \(\iota\) is nonsquare. For every such prime \(p\ge29\),
Hasse's lower bound on the product cover exceeds 16 and therefore its
boundary count, giving the required unit point.

If \(p\equiv1\pmod8\), the base has at most sixteen geometric
zero-coordinate points, giving at most 32 points on the double cover.
Together with its sixteen points above \(b=0\), at most 48 points
are excluded. For \(p\ge73\), Hasse's lower bound exceeds 48.
Below 73 the primes in this class are 17 and 41. At 41 take
\(\iota=9\) and \((V_0,V_1,V_2,V_3)=(6,18,2,13)\).
The squared coordinates are \((36,37,4,5)\), so \(b=1\),
\(Q=32\), and \(2Q=23=8^2\pmod{41}\).
Swapping \(V_1,V_2\) preserves \(f\) and handles the other
choice of \(\iota\). The exclusions follow from
Theorems~\ref{thm:crit-basic} and
\ref{thm:crit-17}, including every root sign.

For local realization repeat
\eqref{crit:lift-coordinates}--\eqref{crit:lift-denominator} with
the prescribed signed pattern. Hensel's lemma gives
\(A_j=p^sv_j\), \(v_j\equiv V_j\pmod p\). The product is
\[
 t=\frac{v_0v_3+v_1v_2}{p^{2s}M^2}.
\]
Twice the numerator is a unit square, so \(2t\) is square.
The ratio \((t+1)/t\in1+p^{2s}\Z_p\) supplies the second
root. Conversely, reduction gives an admissible signed pattern.
\end{proof}

\subsection{Excess radius valuations}

\begin{corollary}\label{cor:crit-excess}
Suppose at least one of \(t,2t,t+1,2(t+1)\) is a
rational square. If \(p\mid G\) and \(v_p(A_j)>v_p(G)\) for
some corner, then
\[
 v_p(L)=2v_p(G),\qquad p\equiv1\pmod{16},\qquad
 2^{(p-1)/4}\equiv1\pmod p.
\]
\end{corollary}

\begin{proof}
At a noncritical prime all normalized radii are units. Hence \(p\)
is critical. A square symmetry moves the unique zero to corner zero
and preserves the opposite-product sum. Dividing by \(a_1\) gives
\[
 a_0=0,\quad a_1=1,\quad a_2^2=\iota,\quad
 a_3^2=1+\iota,\quad Q=a_2.
\]
Thus \(p\equiv1\pmod8\), and 2 is square in \(\Q_p\).
By \eqref{crit:product-class}, any of the four square hypotheses
forces \(\chi_p(Q)=1\). But
\[
 \chi_p(a_2)=\iota^{(p-1)/4}=(-1)^{(p-1)/8},
\]
so \(p\equiv1\pmod{16}\). Finally \((1+\iota)^2=2\iota\) gives
\[
 1=\chi_p(1+\iota)
 =2^{(p-1)/4}(-1)^{(p-1)/8}=2^{(p-1)/4}.
\]
\end{proof}
